\documentclass[11pt]{article}

\usepackage[margin=1.0in]{geometry}
\usepackage{lecnotes}
\input{lmacros}

\newcommand{\course}{15-316: Software Foundations of Security \& Privacy}
\newcommand{\lecturer}{}
\newcommand{\lecdate}{Due Tue, Mar 17, 2026}
\newcommand{\lecnum}{4}
\newcommand{\lectitle}{Homework 4}
\newcommand{\courseurl}{https://15316-cmu.github.io/2026/}

\title{Assignment \lecnum \\ Information Flow}
\date{\lecdate \\ 90 points + 20 points extra credit}

\lhead[\fancyplain{}{\bfseries A\lecnum.\thepage}]%
      {\fancyplain{}{\bfseries\lectitle}}
\chead[]{}
\rhead[\fancyplain{}{\bfseries\lectitle}]%
      {\fancyplain{}{\bfseries A\lecnum.\thepage}}
\lfoot[{\small\scshape Assignments}]{{\small\scshape Assignments}}
\cfoot[]{}
\rfoot[{\small\scshape\lecdate}]{{\small\scshape\lecdate}}

\begin{document}

\maketitle

Your solution should be handed in as a file \texttt{hw4.pdf} to Gradescope.  If
at all possible, write your solutions in \LaTeX.  The handout
\texttt{hw4-infoflow.zip} includes the \LaTeX\ sources for Lectures 12 and 13 and
the necessary style files which provide some examples for rules, derivations,
and proofs.  You may use up to two late days as usual.

\section{Implicit Flows [60 points + 20 points extra credit]}
\label{sec:implicit}

\newcommand{\trycatch}{\mb{try/catch}}
\newcommand{\try}[2]{\mb{try}\; #1\; \mb{catch}\; #2}

Consider adding a new construct to our language \textsc{SafeTiny},
$\mb{try}\; \alpha\; \mb{catch}\; \beta$.  Note that in \textsc{SafeTiny} all
commands are safe, but we have $\mb{test}\; P$ which aborts if $P$ is false.  We
do not consider division, memory read/write, or $\mb{assert}$.  In order to
simplify matters further, we also exclude loops from consideration, but see
the extra credit tasks at the end of this problem.

$\try{\alpha}{\beta}$ is supposed to execute as follows:
\begin{enumerate}
\item Execute $\alpha$ in the current state $\omega$
\item If $\alpha$ \textbf{does not abort}, the $\mb{try}/\mb{catch}$ construct
  finishes in the poststate of $\alpha$
\item If $\alpha$ \textbf{aborts}, we continue by executing $\beta$ in the
  prestate $\omega$.  In this case, the poststate of $\beta$ will be the
  poststate of $\try{\alpha}{\beta}$.
\end{enumerate}
$\trycatch$ is not easy to implement efficiently in a compiler since we have to
either save the prestate $\omega$, or track assignments so we can roll back the
state when a test fails.  As we see in \autoref{task:eval}, it is not so
difficult in an interpreter.

Here are some examples:
\[
  \begin{array}{lcl}
    \eval\; \omega\; (\try{\mb{test}\; \bot}{x := 0})
    & = & \omega[x \mapsto 0] \\
    \eval\; \omega\; (\try{\mb{test}\; \top}{x := 0})
    & = & \omega \\
    \eval\; \omega\; (\try{\mb{test}\; \bot}{\mb{test}\; \bot})
    & & \mbox{aborts} \\
    \eval\; \omega\; (\try{(x := 0 \semi \try{(\mb{test}\; x > 0)}{x := 1})}{x := 2})
    & = & \omega[x \mapsto 1]
  \end{array}
\]

\begin{task}[10 points]\rm
  Give a semantic definition of $\omega \lbb \try{\alpha}{\beta} \rbb \nu$ and
  $\mb{test}\; P$ that models the intended behavior based on the informal
  description above.  You should model a program that aborts (and is not caught)
  as one that has no poststate.
\end{task}

With this understanding, we can update our definition of $\eval$ so that it
explicitly returns either a state $\nu$ or $\ms{abort}$.  We show the cases for
sequential composition, $\mb{skip}$, and assignment.
\[
  \begin{array}{lcll}
    \eval\; \omega\; (\alpha \semi \beta) & = & \ms{abort}
    & \mbox{if $\eval\; \omega\; \alpha = \ms{abort}$} \\
    \eval\; \omega\; (\alpha \semi \beta) & = & \ms{abort}
    & \mbox{if $\eval\; \omega\; \alpha = \mu$ and $\eval\; \mu\; \beta = \ms{abort}$} \\
    \eval\; \omega\; (\alpha \semi \beta) & = & \nu
    & \mbox{if $\eval\; \omega\; \alpha = \mu$ and $\eval\; \mu\; \beta = \nu$} \\[1ex]
    \eval\; \omega\; (\mb{skip}) & = & \omega \\
    \eval\; \omega\; (x := e) & = & \omega[x \mapsto c]
    & \mbox{where $\evalz\; \omega\; e = c$}
  \end{array}
\]

\begin{task}[15 points]\rm
  \label{task:eval}
  Complete the definition of $\eval$ with the cases for conditionals,
  tests, and $\trycatch$.
\end{task}

\begin{task}[10 points]\rm
  Conjecture an axiom of equivalence for $[\try{\alpha}{\beta}]Q$, or explain
  briefly why you believe no such axiom is possible in dynamic logic (as we have
  constructed it so far).  Note that your axiom only needs to be sound in the
  language without loops.
\end{task}

\begin{task}[5 points]\rm
  Give an example of a security policy $\Sigma_0$ and program $\alpha_0$
  demonstrating that $\mb{test}$ and $\trycatch$ create a new possibility for
  implicit information flow.  For this question, you should work with a security
  lattice with just two elements $\ms{H}$ and $\ms{L}$ with
  $\ms{L} \sqsubset \ms{H}$ and the definition of termination-insensitive
  noninterference from \href{\courseurl/lectures/12-infoflow.pdf}{Lecture 12}.
\end{task}

\begin{task}[10 points]\rm
  Prove that your example from the previous task violates
  termination-insensitive noninterference, that is,
  $\Sigma_0 \models \alpha_0\; \ms{secure}$ is \textbf{not} true.
\end{task}

In order to prevent the implicit flows enabled by $\trycatch$ we introduce a new
ghost variable $\mi{handler}$ into the information flow type system.  The
security level of $\mi{handler}$ should be that of the $\mb{catch}$ that would
be invoked should the current program abort.  It should be $\bot$ (the least
element of the security lattice) at the beginning of evaluation.

\begin{task}[10 points]\rm
  Give rules for $\trycatch$ and $\mb{test}$ in the information flow type
  system.  You do not need to prove their soundness.
\end{task}

\begin{quote}\it
  \color{red}The remainder of Problem 1 is for extra credit.
\end{quote}

\noindent
In order to support loops, we assume a global bound $\ms{b}$ on the number of
iterations for each while loop, after which it aborts.  For example, with
$\ms{b} = 0$ the program aborts if it ever attempts to enter the body of a loop,
with $\ms{b} = 1$ each loop $\mb{while}\; P\; \alpha$ is equivalent to
$\mb{if}\; P\; \mb{then}\; (\alpha \semi \mb{test}\; (\lnot P)) \; \mb{else}\;
\mb{skip}$.

\begin{task}[\color{red}5 bonus points]\rm
  Give a semantic definition of $\omega \lbb \mb{while}\; P\; \alpha\rbb \nu$
  that models the intended behavior of bounded loops based on the informal
  description above.
\end{task}

\begin{task}[\color{red}5 bonus points]\rm
  Complete the definition $\eval$ by providing a clause for $\mb{while}$ loops
  bounded by $\ms{b}$.  Feel free to use auxiliary functions.
\end{task}

In order the conjecture suitable axioms in dynamic logic assume that each loop
$\mb{while}\; P\; \alpha$ with loop invariant $J$ is written explicitly as
$\mb{while}_J^{\ms{b}}\; P\; \alpha$.

\begin{task}[\color{red}10 bonus points]\rm
  Conjecture axioms in dynamic logic for $\mb{while}^{\ms{b}}_J$ and
  $\trycatch$, or explain why you think that bounded loops and $\trycatch$
  cannot be axiomatized in the framework of dynamic logic (as we have
  constructed it so far).
\end{task}

\section{Declassification [30 points]}

Consider the formulation of termination-insensitive noninterference in the
presence of declassification under the two-level security lattice
($\ms{L} \sqsubset \ms{H}$).

\begin{quote}
  Assume $\alpha$ contains a single occurrence of $\mb{declassify}_{\ms{L}}(e)$
  where $x \in \ms{use}\; e$ implies $x \not\in \ms{maydef}\; \alpha$.
  \par \medskip For such programs we define
  $\Sigma \models \alpha\; \ms{secure}$ iff \newline whenever
  $\Sigma \vdash \omega_1 \approx_{\ms{L}} \omega_2$ \newline and
  $\eval\; \omega_1\; e = \eval\; \omega_2\; e$ \newline then
  $\Sigma \vdash \eval\; \omega_1\; \alpha \approx_{\ms{L}} \eval\; \omega_2\;
  \alpha$
\end{quote}

\begin{task}[20 points]\rm
  Assume you are given a security policy $\Sigma$ and a program $\alpha$ that
  contains a single occurrence of $\mb{declassify}_{\ms{L}}(e)$ satisfying the
  use/maydef condition.  Show how to construct a formula $R$ in dynamic logic
  such that the validity of $R$ implies $\Sigma \models \alpha\; \ms{secure}$.
  Your starting point should be the construction in Section 2 of
  \href{\courseurl/lectures/13-declassification.pdf}{Lecture 13}.
\end{task}

\begin{task}[5 points]\rm
  Give an example policy $\Sigma_0$ and program $\alpha_0$ that is \textbf{not}
  secure due to incorrect use of declassification, that is, the use/maydef
  condition is violated.  Show the encoding of the example in dynamic logic and
  determine whether it is valid.  Explain your finding.
\end{task}

\begin{task}[5 points]\rm
  Give an example policy $\Sigma_1$ and a program $\alpha_1$ that uses
  declassification and is secure.  Show the encoding of the example in dynamic
  logic and demonstrate that it is valid.  This is usually done most directly by
  constructing a weakest precondition, if the formula is in the class that
  permits it.  You don't need to show intermediate steps.
\end{task}

\end{document}
